%% ============================================================================
%% Reviewer point addressed
%% ----------------------------------------------------------------------------
%% "Reproducibility needs improvement. The authors should provide details on
%%  random seeds, instance-generation procedure, training/test splits, exact
%%  sampling of disclosure times, parameter counts, GPU memory use, and whether
%%  code/data will be released. This is especially important because the
%%  baseline implementations are modified rather than used directly."
%%
%% Every number below is measured by reproduction/evaluate_table1.py and
%% reproduction/check_protocol.py on an NVIDIA RTX 4070 SUPER; regenerate with
%%   python -u -m reproduction.check_protocol
%%   python -u -m reproduction.evaluate_table1 --measure-memory
%% ============================================================================


%% ============================================================================
%% PATCH 1 -- Replace the "Node Revelation Process." paragraph of Section IV-A
%% so that the disclosure process is stated exactly rather than as a scalar rate.
%% ============================================================================
%% 【REPLACE】the paragraph beginning "Node Revelation Process. Node revelations
%%            are assumed to follow a Poisson distribution..." together with
%%            Eq. (34) and the sentence defining lambda.
%% 【WITH】:

\textbf{Node Revelation Process.} Disclosure times are drawn from a Poisson
process discretized over the ordered customer slots. Static customers receive
disclosure time $0$. A dynamic customer occupying ordered slot
$i \in \{0,\ldots,n-1\}$ receives
\begin{equation}
a_i \;=\; \mathrm{clip}\!\left(\mathrm{Poisson}(\lambda_i),\, 1,\, T\right),
\qquad
\lambda_i \;=\; 1 + \frac{i\,(T-1)}{n-1},
\label{eq:disclosure}
\end{equation}
so that the per-slot rates are spaced linearly from $1$ to $T$ minutes and every
draw is independent across instances. Their mean is exactly
$\bar{\lambda} = (1+T)/2 = 240.5$ minutes, which is the scalar rate reported in
earlier drafts; we state the per-slot rates here because they, and not the
scalar mean, define the process that is sampled. For $n=20$ the realized mean
disclosure times of dynamic customers on the test split are $256.8$, $248.7$,
$238.7$, and $240.8$ minutes at $\phi = 10\%$, $25\%$, $50\%$, and $75\%$.

A disclosure time is distinguished from the time at which a planner may act on
it. The neural methods replan on the interval clock, so a request disclosed at
time $a_i$ first enters the available set at the earliest interval boundary
$T_r \geq a_i$, as specified in Algorithm~\ref{algo:interval_transition}. The
Greedy heuristic is event-driven and reacts whenever a vehicle becomes idle, so
it acts at $a_i$ itself.


%% ============================================================================
%% PATCH 2 -- Extend the "Training Setup." paragraph of Section IV-A with the
%% instance-generation procedure, the splits, and the seeds.
%% ============================================================================
%% 【APPEND】 at the end of the "Training Setup." paragraph:

\textbf{Instances and splits.} Each instance draws the depot and $n$ customer
locations independently and uniformly from the unit square, demands from the
discrete uniform distribution on $\{5,\ldots,41\}$, and service durations from
the discrete uniform distribution on $\{10,\ldots,31\}$ minutes. Exactly
$\mathrm{round}(n\phi)$ customers are dynamic, selected uniformly without
replacement, and receive disclosure times from~\eqref{eq:disclosure}. Vehicle
capacity is consumed over the whole horizon and is not replenished, which is the
single-round-trip premise used by our QoS definition. Coordinates remain in
$[0,1]$; demands are divided by the capacity and all temporal features by $T$,
so the normalized horizon is $1$ and the normalized speed is $T$.

The four dynamic rates are evaluated on paired instances: coordinates, demands,
service durations, and disclosure draws are shared across rates and the dynamic
sets are nested, so increasing $\phi$ only converts further requests from static
to dynamic. Under this common-random-numbers design a cost difference between
two rates reflects the dynamic rate rather than a difference in instance
difficulty.

Training, validation, and test instances are drawn from three disjoint streams
with fixed seeds $1234$, $4321$, and $20260821$; evaluation decoding uses seed
$314159$. The validation split contains $100$ instances per dynamic rate and is
used only to select the reported checkpoint, by validation QoS first and then
validation executed-route cost. The test split contains $100$ instances per
dynamic rate and is evaluated once, after the checkpoint and the decoding rule
are frozen. Because the instances are a deterministic function of the recorded
seeds, the splits can be regenerated exactly; we nonetheless release the
tensors together with SHA-256 hashes so that a reader can confirm that the split
they regenerate is the one behind Table~\ref{tab:synthetic}.


%% ============================================================================
%% PATCH 3 -- Replace the "Baselines." paragraph's opening so the nature of the
%% baseline adaptation, and its cost, are stated plainly.
%% ============================================================================
%% 【REPLACE】: "Since no existing NCO method directly addresses DCVRP under our
%%              problem formulation, we integrate only the vehicle selection
%%              strategy of each baseline into our unified time-driven
%%              framework, while keeping all other components (encoder, node
%%              decoder, training procedure, and dynamic environment)
%%              identical. This controlled setup ensures a fair comparison by
%%              isolating the effect of the vehicle coordination mechanism."
%% 【WITH】:

Since no existing NCO method addresses DCVRP under our problem formulation, none
of the baselines can be executed unmodified on these instances. We therefore
re-implement the \emph{vehicle-selection strategy} of each baseline inside our
time-driven environment and hold every other component identical: the encoder,
the node decoder, the training procedure, the reward, the instance
distribution, and the dynamic environment are shared. The shared
encoder--decoder accounts for $743{,}040$ parameters in all five methods, and
the only architectural difference between them is the selection module, so the
comparison isolates the vehicle-coordination mechanism. We state the
consequence explicitly: these rows measure the published \emph{coordination
strategies} transplanted into a common environment, not the baselines' original
released systems, and they should not be read as a reproduction of the numbers
reported in the corresponding papers.


%% ============================================================================
%% PATCH 4 -- New paragraph for the "Model Setup." part of Section IV-A giving
%% parameter counts and memory. Insert after "...rollout baseline with 3
%% rollouts and a threshold set to 0.05."
%% ============================================================================
%% 【INSERT】:

\textbf{Parameter counts and memory.} Table~\ref{tab:budget} reports the
parameter budget of each method, separated into the shared encoder--decoder and
the vehicle-selection module, together with peak GPU memory. A training step
comprises the sampled policy rollout, the three baseline rollouts, the backward
pass, and the optimizer step at batch size $100$. MARDAM's earliest-idle rule
and MAAM's round-robin rule contain no trainable selection parameters, which is
a property of the published strategies rather than of our implementation: for
those two methods the shared node decoder is the only component that learns.
Although the reported experiments were run on a 4090D with 24~GB, the $n=20$
setting requires under $2$~GB for training and under $0.1$~GB for inference, so
it reproduces on commodity hardware; the measurements in
Table~\ref{tab:budget} were taken on an RTX~4070~SUPER with 12~GB.

\begin{table}[t]
\centering
\caption{Parameter budget and peak GPU memory at $n=20$, $m=4$. The shared
encoder--decoder is identical across methods; only the vehicle-selection module
differs.}
\label{tab:budget}
\begin{tabular}{lrrrrr}
\toprule
& \multicolumn{3}{c}{Parameters} & \multicolumn{2}{c}{Peak GPU memory (GiB)} \\
\cmidrule(lr){2-4}\cmidrule(lr){5-6}
Method & Shared & Selector & Total & Training & Inference \\
\midrule
MARDAM & 743{,}040 & 0       & 743{,}040 & 0.95 & 0.03 \\
MAAM   & 743{,}040 & 0       & 743{,}040 & 0.96 & 0.05 \\
LiDRL  & 743{,}040 & 150{,}404 & 893{,}444 & 1.22 & 0.05 \\
AMCVN  & 743{,}040 & 83{,}841  & 826{,}881 & 1.14 & 0.05 \\
DVNDA  & 743{,}040 & 183{,}556 & 926{,}596 & 1.76 & 0.05 \\
\bottomrule
\end{tabular}
\end{table}


%% ============================================================================
%% PATCH 5 -- New paragraph closing Section IV-A: the training-free protocol
%% check. This is the strongest single reproducibility claim available, because
%% Greedy contains no learned component.
%% ============================================================================
%% 【INSERT】 at the end of Section IV-A:

\textbf{Protocol verification.} The Greedy row of
Table~\ref{tab:synthetic} contains no learned component, so it verifies the
experimental protocol independently of model quality: any disagreement is
attributable to the instance distribution, the time and capacity semantics, or
the cost definition rather than to training. Re-running the released
reproduction package regenerates that row to within $1.8\%$ at every dynamic
rate ($9.03$, $9.71$, $11.31$, and $12.21$ against the reported $9.07$, $9.69$,
$11.25$, and $12.43$), with QoS between $99.70\%$ and $99.80\%$. We report this
check because it lets a reader confirm the environment before spending any GPU
time on training.


%% ============================================================================
%% PATCH 6 -- Code and data availability. Replace the footnote/abstract sentence
%% "Our code is available at https://github.com/ftwangyang/DCVRP."
%% ============================================================================
%% 【REPLACE】 with:

Our code, trained weights, the fixed evaluation manifests with SHA-256 hashes,
and the raw per-instance results are available at
\url{https://github.com/ftwangyang/DCVRP}. The release includes the immutable
experimental protocol, the instance generator, the training and evaluation
entry points, and a training-free script that verifies the protocol before any
model is trained.


%% ============================================================================
%% PATCH 7 (optional) -- If the reviewer's concern about released-code fidelity
%% is raised again, this paragraph can be added to the supplementary material.
%% ============================================================================
%% 【SUPPLEMENTARY】:

\textbf{Differences between the released implementation and the description
above.} Three details of the original implementation differ from the
specification we now state, and we record them so that a reader comparing the
two is not misled. First, the released generator draws integer coordinates in
$[0,100]$ and min--max normalizes them over the generated batch, and it assigns
dynamic status by independent Bernoulli$(\phi)$ draws, so an instance labelled
$25\%$ need not contain exactly five dynamic customers; the specification above
uses continuous coordinates and exact dynamic counts. Second, the released
generator draws one vector of $n$ Poisson values per batch rather than per
instance, so instances within a batch share disclosure times; the specification
above draws independently per instance. Third, the released environment retains
$T=480$ for its interval boundaries after the temporal features have already
been divided by $T$, which makes every dynamic request visible at the first
boundary; under that reading Greedy's cost is non-monotonic in $\phi$
($10.24$, $12.84$, $13.99$, $12.61$), peaking at $\phi=50\%$, whereas
Table~\ref{tab:synthetic} increases monotonically. The unit-horizon convention
stated above is the one consistent with both the normalization and
Table~\ref{tab:synthetic}. The reproduction package implements all variants and
reports the measured difference between them.
